% Blueprint for: Global Steady State of the Vlasov-Maxwell-Landau System
% Based on Guo & Strain, "Almost exponential decay near Maxwellian"

\chapter{Introduction}

This blueprint describes the formalization of the characterization of global steady states of the Vlasov-Maxwell-Landau (VML) system on a flat 3-torus $\T^3$. The main result (Theorem~42 in the notation of the original paper) states that any smooth, positive, integrable steady state of the VML system with sufficient velocity decay must be a spatially homogeneous equilibrium Maxwellian, with vanishing electric field and constant magnetic field.

\chapter{Definitions}\label{ch:defs}

\begin{definition}[Norm squared]\label{def:normSq}
  \lean{VML.normSq}
  \leanok
  For $v \in \R^3$, define $|v|^2 = \sum_{i=1}^{3} v_i^2$.
\end{definition}

\begin{definition}[Euclidean norm]\label{def:eucNorm}
  \lean{VML.eucNorm}
  \leanok
  \uses{def:normSq}
  For $v \in \R^3$, define $|v| = \sqrt{|v|^2}$.
\end{definition}

\begin{definition}[Landau matrix]\label{def:landauMatrix}
  \lean{VML.landauMatrix}
  \leanok
  \uses{def:normSq, def:eucNorm}
  Given a kernel function $\Psi : \R \to \R$, the Landau matrix $A(z)$ is the $3 \times 3$ matrix with entries
  \[
    A_{ij}(z) = \Psi(|z|)\bigl(\delta_{ij} |z|^2 - z_i z_j\bigr).
  \]
\end{definition}

\begin{definition}[Maxwellian]\label{def:isMaxwellian}
  \lean{VML.IsMaxwellian}
  \leanok
  A function $f : \R^3 \to \R$ is \emph{Maxwellian} if there exist $a_0 \in \R$, $b \in \R^3$, and $c_0 < 0$ such that
  \[
    f(v) = \exp\bigl(a_0 + b \cdot v + c_0 |v|^2\bigr)
  \]
  for all $v \in \R^3$.
\end{definition}

\begin{definition}[Equilibrium Maxwellian]\label{def:equilibriumMaxwellian}
  \lean{VML.equilibriumMaxwellian}
  \leanok
  For $\rho_{\mathrm{ion}} > 0$ and $T > 0$, the equilibrium Maxwellian is
  \[
    M_{\rho,T}(v) = \frac{\rho_{\mathrm{ion}}}{(2\pi T)^{3/2}} \exp\!\left(-\frac{|v|^2}{2T}\right).
  \]
\end{definition}

\begin{definition}[Landau collision operator]\label{def:landauOperator}
  \lean{VML.LandauOperator}
  \leanok
  \uses{def:landauMatrix}
  The Landau collision operator is
  \[
    Q(f, f)(v) = \grad_v \cdot \int_{\R^3} A(v - w)\bigl[f(w)\,\grad_v f(v) - f(v)\,\grad_w f(w)\bigr]\, dw.
  \]
\end{definition}

\begin{definition}[Entropy dissipation]\label{def:entropyDissipation}
  \lean{VML.entropyDissipation}
  \leanok
  \uses{def:landauOperator}
  The entropy dissipation functional is
  \[
    D(f) = \int_{\R^3} Q(f, f)(v)\, \log f(v)\, dv.
  \]
\end{definition}

\begin{definition}[Flat 3-torus]\label{def:flatTorus3}
  \lean{VML.FlatTorus3}
  \leanok
  A \emph{flat 3-torus} is a compact, nonempty, first-countable measure space $X$ equipped with spatial gradient, divergence, and curl operators, a differentiability predicate \texttt{IsSpatiallyDiff}, and axioms encoding:
  \begin{itemize}
    \item operator linearity (gradient of constants, sums, scalar multiples, chain rules),
    \item closure of \texttt{IsSpatiallyDiff} under const, add, smul, log, grad, and continuity,
    \item integration by parts ($\int \varphi\, \partial_i \psi = -\int \psi\, \partial_i \varphi$),
    \item curl-integral vanishing ($\int u \cdot \curl F = 0$),
    \item harmonic functions are constant,
    \item maximum principle for the Laplacian,
    \item Killing fields are harmonic; curl-free divergence-free fields are harmonic,
    \item Fubini theorem for spatial--velocity double integrals.
  \end{itemize}
\end{definition}

\begin{definition}[VML steady-state input]\label{def:vmlInput}
  \lean{VML.VMLInput}
  \leanok
  \uses{def:flatTorus3, def:landauOperator, def:entropyDissipation, def:isMaxwellian}
  The \texttt{VMLInput} structure bundles the physical data $(f, E, B, \Psi, \nu, \rho_{\mathrm{ion}})$ together with:
  positivity of $f$, $\nu$, $\rho_{\mathrm{ion}}$, $\Psi$;
  smoothness in $v$; integrability;
  the steady-state Maxwell equations (Amp\`ere, Gauss, $\divop B = 0$);
  the steady-state Vlasov equation;
  spatial differentiability of $f(\cdot, v)$ and $B$;
  and all derived analytical results needed to close the argument (entropy dissipation vanishing, Maxwellian form, polynomial identity, Poisson--Boltzmann, etc.).
\end{definition}

\begin{definition}[Velocity decay conditions]\label{def:velocityDecayConditions}
  \lean{VML.VelocityDecayConditions}
  \leanok
  The \texttt{VelocityDecayConditions} structure bundles 17 integrability/Fubini/IBP conditions on $f$, $E$, $B$ that are needed to justify the analytical manipulations (integration by parts on $\R^3$, Fubini exchange, Leibniz differentiation under the integral, entropy dissipation continuity, etc.).
\end{definition}


\chapter{Landau Matrix Properties}\label{ch:section2}

\begin{lemma}[Positive semi-definiteness]\label{lem:landauMatrix_posSemidef}
  \lean{VML.landauMatrix_posSemidef}
  \leanok
  \uses{def:landauMatrix}
  If $\Psi(|z|) \ge 0$, then the Landau matrix $A(z)$ is positive semi-definite:
  \[
    Y^T A(z) Y \ge 0 \quad \text{for all } Y \in \R^3.
  \]
\end{lemma}

\begin{lemma}[Nullspace characterization]\label{lem:landauMatrix_quadForm_eq_zero_iff}
  \lean{VML.landauMatrix_quadForm_eq_zero_iff}
  \leanok
  \uses{def:landauMatrix}
  If $\Psi(|z|) > 0$ and $z \ne 0$, then $Y^T A(z) Y = 0$ if and only if $Y$ is parallel to $z$ (i.e., $\exists \lambda,\, Y = \lambda z$).
\end{lemma}


\chapter{H-Theorem and Maxwellian Characterization}\label{ch:section3}

\begin{theorem}[H-theorem]\label{thm:H_theorem}
  \lean{VML.H_theorem}
  \leanok
  \uses{def:entropyDissipation, lem:landauMatrix_posSemidef}
  If $\Psi \ge 0$ and $f > 0$ is smooth and integrable, then the entropy dissipation is non-positive:
  \[
    D(f) \le 0.
  \]
\end{theorem}

\begin{lemma}[Logarithm is quadratic]\label{lem:log_f_quadratic}
  \lean{VML.log_f_quadratic}
  \leanok
  If $f > 0$, $f$ smooth, and $\grad_v \log f(v) = b + 2c_0 v$ for all $v$, then there exists $a_0$ such that $\log f(v) = a_0 + b \cdot v + c_0 |v|^2$.
\end{lemma}

\begin{theorem}[Nullspace necessity]\label{thm:nullspace_necessity}
  \lean{VML.nullspace_necessity}
  \leanok
  \uses{def:isMaxwellian, def:entropyDissipation, lem:landauMatrix_quadForm_eq_zero_iff, lem:log_f_quadratic}
  If $\Psi > 0$, $f > 0$, $f$ smooth and integrable, $Q(f,f) = 0$, and the PSD/score-form conditions hold, then $f$ is Maxwellian.
\end{theorem}

\begin{theorem}[Nullspace sufficiency]\label{thm:nullspace_sufficiency}
  \lean{VML.nullspace_sufficiency}
  \leanok
  \uses{def:landauOperator, def:isMaxwellian}
  If $\log f(v) = a_0 + b \cdot v + c_0 |v|^2$ with $c_0 < 0$, then $Q(f,f)(v) = 0$ for all $v$.
\end{theorem}


\chapter{Transport Constraints}\label{ch:section4}

\begin{theorem}[Steady state is local Maxwellian]\label{thm:steady_state_is_local_maxwellian}
  \lean{VML.steady_state_is_local_maxwellian}
  \leanok
  \uses{thm:nullspace_necessity}
  If $\Psi > 0$, $f > 0$, $f$ is smooth, and $D(f(x, \cdot)) = 0$ for all $x$, then for every $x$, $f(x, \cdot)$ is Maxwellian.
\end{theorem}

\begin{lemma}[Lorentz force is divergence-free]\label{lem:lorentz_force_div_zero}
  \lean{VML.lorentz_force_div_zero}
  \leanok
  For any $E_0, B_0 \in \R^3$,
  \[
    \divop_v \bigl(E_0 + v \times B_0\bigr) = 0.
  \]
\end{lemma}

\begin{theorem}[Transport entropy identity]\label{thm:transport_entropy_from_vlasov}
  \lean{VML.transport_entropy_from_vlasov}
  \leanok
  \uses{def:flatTorus3, def:entropyDissipation}
  Under the Vlasov equation $v \cdot \grad_x f + (E + v \times B) \cdot \grad_v f = \nu\, Q(f,f)$, combined with integration by parts (spatial and velocity) and the divergence-free property of the Lorentz force, one derives
  \[
    \int_X D(f(x, \cdot))\, dx = 0.
  \]
\end{theorem}


\chapter{Polynomial Matching}\label{ch:section5}

\begin{lemma}[Temperature is constant]\label{lem:temperature_constant}
  \lean{VML.temperature_constant}
  \leanok
  \uses{def:flatTorus3}
  If for all $x$ and $v$,
  $(v \cdot \grad_x c(x))\, |v|^2 = 0$,
  then $\grad_x c(x) = 0$ for all $x$.
\end{lemma}

\begin{theorem}[Polynomial identity from Vlasov]\label{thm:polynomial_identity_from_vlasov}
  \lean{VML.polynomial_identity_from_vlasov}
  \leanok
  \uses{def:flatTorus3, def:isMaxwellian}
  If $f$ is a local Maxwellian $f(x, v) = \exp(a(x) + b(x) \cdot v + c(x)|v|^2)$ satisfying the Vlasov equation, then matching coefficients of the polynomial in $v$ yields:
  $\grad c = 0$ (constant temperature),
  and transport identities for the lower-order coefficients.
\end{theorem}


\chapter{Bulk Velocity}\label{ch:section6}

\begin{theorem}[Bulk velocity vanishes]\label{thm:bulk_velocity_zero}
  \lean{VML.bulk_velocity_zero}
  \leanok
  \uses{def:flatTorus3, thm:polynomial_identity_from_vlasov}
  In steady state with constant temperature and Amp\`ere's law, the bulk velocity is zero: $b_0 = 0$.
\end{theorem}


\chapter{Electric Field}\label{ch:section7}

\begin{theorem}[Poisson--Boltzmann from Vlasov]\label{thm:poisson_boltzmann_from_vlasov}
  \lean{VML.poisson_boltzmann_from_vlasov}
  \leanok
  \uses{def:flatTorus3, thm:bulk_velocity_zero}
  When $b_0 = 0$ (isotropic Maxwellian), combining the Vlasov equation with Gauss's law yields the Poisson--Boltzmann equation:
  \[
    T_\infty \,\Delta(\log \rho) = \rho - \rho_{\mathrm{ion}}.
  \]
\end{theorem}

\begin{theorem}[Electric field vanishes]\label{thm:electric_field_zero}
  \lean{VML.electric_field_zero}
  \leanok
  \uses{thm:poisson_boltzmann_from_vlasov}
  By the maximum principle on the compact torus (applied to $\log \rho$), the density $\rho$ must be constant, and therefore $E(x) = 0$ for all $x$.
\end{theorem}


\chapter{Magnetic Field}\label{ch:section8}

\begin{theorem}[Magnetic field is constant]\label{thm:magnetic_field_constant}
  \lean{VML.magnetic_field_constant}
  \leanok
  \uses{def:flatTorus3, thm:bulk_velocity_zero}
  With $J = 0$ (since $b_0 = 0$), Amp\`ere's law gives $\curl B = 0$, and combined with $\divop B = 0$ and the fact that curl-free divergence-free vector fields on the torus are harmonic (hence constant), we conclude: there exists $B_0$ such that $B(x) = B_0$ for all $x$.
\end{theorem}


\chapter{Main Results}\label{ch:main}

\begin{theorem}[Main from physics]\label{thm:main_from_physics}
  \lean{VML.main_from_physics}
  \leanok
  \uses{def:vmlInput, thm:steady_state_is_local_maxwellian, thm:polynomial_identity_from_vlasov, thm:bulk_velocity_zero, thm:electric_field_zero, thm:magnetic_field_constant, def:equilibriumMaxwellian}
  Given a \texttt{VMLInput}, there exist $T_{\mathrm{eq}} > 0$ and $B_0 \in \R^3$ such that:
  \begin{enumerate}
    \item $f(x, v) = M_{\rho_{\mathrm{ion}}, T_{\mathrm{eq}}}(v)$ for all $x, v$,
    \item $E(x) = 0$ for all $x$,
    \item $B(x) = B_0$ for all $x$.
  \end{enumerate}
\end{theorem}

\begin{theorem}[Theorem 42]\label{thm:theorem42}
  \lean{VML.Theorem42}
  \leanok
  \uses{thm:main_from_physics, def:velocityDecayConditions, def:flatTorus3, def:equilibriumMaxwellian}
  Let $X$ be a flat 3-torus. Let $f : X \times \R^3 \to \R$ be a smooth, positive, integrable function, let $E, B : X \to \R^3$, and let $\Psi$ be a positive collision kernel. Suppose:
  \begin{itemize}
    \item The steady-state Vlasov equation holds: $v \cdot \grad_x f + (E + v \times B) \cdot \grad_v f = \nu\, Q_\Psi(f, f)$.
    \item The steady-state Maxwell equations hold (Amp\`ere, Gauss, $\divop B = 0$).
    \item $f(\cdot, v)$ is spatially differentiable for each $v$, and $B$ is spatially differentiable.
    \item The velocity decay conditions \texttt{VelocityDecayConditions} are satisfied.
  \end{itemize}
  Then there exist $T_{\mathrm{eq}} > 0$ and $B_0 \in \R^3$ such that:
  \[
    f(x, v) = \frac{\rho_{\mathrm{ion}}}{(2\pi T_{\mathrm{eq}})^{3/2}} \exp\!\left(-\frac{|v|^2}{2T_{\mathrm{eq}}}\right), \quad
    E(x) = 0, \quad B(x) = B_0
  \]
  for all $x \in X$ and $v \in \R^3$.
\end{theorem}

\begin{theorem}[Coulomb concrete Theorem 42]\label{thm:coulomb_concrete}
  \lean{VML.CoulombConcreteTheorem42}
  \leanok
  \uses{thm:theorem42}
  Specialization of Theorem~42 to the Coulomb kernel $\Psi(r) = r^{-3}$ on the concrete torus $\T^3 = (\R / \mathbb{Z})^3$, with Schwartz-class decay assumptions on $f$.
  The conclusion is the same: $f$ is an equilibrium Maxwellian, $E = 0$, and $B$ is constant.
\end{theorem}

\begin{theorem}[Temperature uniqueness]\label{thm:equilibriumMaxwellian_T_unique}
  \lean{VML.equilibriumMaxwellian_T_unique}
  \leanok
  \uses{def:equilibriumMaxwellian}
  If $\rho > 0$, $T_1 > 0$, $T_2 > 0$, and $M_{\rho, T_1} = M_{\rho, T_2}$, then $T_1 = T_2$.
\end{theorem}
